\begin{figure}[htbp]
  \centering
  \begin{tikzpicture}[node distance=12mm and 18mm,>=stealth]
    \node[draw,rounded corners] (as) {几乎处处收敛};
    \node[draw,rounded corners,right=of as] (p) {依概率收敛};
    \node[draw,rounded corners,below=of as] (lp) {$L^p$ 收敛};
    \node[draw,rounded corners,right=of lp] (d) {依分布收敛};
    \draw[->] (as) -- (p);
    \draw[->] (lp) -- (p);
    \draw[->] (p) -- (d);
    \draw[dashed,->] (p) to[bend left=22] node[above,sloped,font=\small] {子列可几乎处处收敛} (as);
  \end{tikzpicture}
  \caption{常用概率收敛关系。虚线表示子列结论，不是原序列的反向蕴含。}
\end{figure}
